\documentclass[tate,book,paper=a5,fontsize=10pt,sidenote_length=6zw]{jlreq}
\begin{document}
\chapter{雨の朝}
\section{駅まで}
朝から細かい雨が降っていた。傘を持たずに家を出たのは失敗だったと、駅までの坂道を歩きながら何度も思った。商店街の軒先を伝って歩けば少しは濡れずに済むが、それでも肩のあたりはすぐに冷たくなった\footnote{この町の商店街には屋根がない。}。

改札の前で友人が待っていた。彼は大きな黒い傘を畳みながら、遅いぞと笑った\marginpar{友人は時間に厳しい。}。約束の時刻にはまだ五分あったが、彼はいつも十分前には着いている。

\section{車内}
電車は思ったより空いていた。窓際の席に座ると、外の景色が雨粒で滲んで見えた。友人は鞄から文庫本を取り出し、しおりを挟んだ頁を開いた。
\begin{itemize}
\item 一つ目の駅で老人が乗ってきた。
\item 二つ目の駅で学生の一団が降りた。
\item 三つ目の駅で車内は静かになった。
\end{itemize}
私は窓の外を眺めながら、昨夜読んだ手紙のことを考えていた。差出人の名前はなく、ただ一行、来週の木曜に例の場所でとだけ書かれていた。

\chapter{木曜日}
\section{例の場所}
例の場所というのが何を指すのか、私にはすぐにわかった。学生の頃によく通った川沿いの喫茶店である。店主は無口な人で、注文を聞く以外にはほとんど口をきかなかった。

木曜日、私は約束の時刻より少し早く店に入った。店内は昔と変わらず薄暗く、奥の席には誰もいなかった。珈琲を頼み、窓の外の川を眺めていると、扉の鈴が鳴った。

入ってきたのは、思いもよらない人物だった。
\end{document}
